\documentclass[aspectratio=169,11pt]{ctexbeamer}

\usetheme{Madrid}
\usecolortheme{default}

\usepackage{amsmath,amssymb,mathtools}
\usepackage{bm}
\usepackage{graphicx}

\makeatletter
\def\input@path{{../模板/}}
\makeatother
\usepackage{yingxing-beamer}

\title{Taylor 展开}
\subtitle{局部多项式逼近、余项估计与常用展开式}
\author{王晨旭}
\date{}

\newcommand{\R}{\mathbb{R}}
\newcommand{\dd}{\,\mathrm{d}}
\newcommand{\ee}{\mathrm{e}}
\newcommand{\abs}[1]{\left|#1\right|}
\newcommand{\ord}{\operatorname{o}}

\begin{document}

\begin{frame}
  \titlepage
\end{frame}

\begin{frame}{本讲目标}
\begin{itemize}
  \item 理解 Taylor 展开是在一点附近用多项式逼近函数。
  \item 区分 Peano 余项、Lagrange 余项、Cauchy 余项和积分余项的用途。
  \item 会判断“展开到几阶”以及余项是否足够小。
  \item 掌握常见函数在 $0$ 处的 Maclaurin 展开及其收敛区间。
\end{itemize}
\end{frame}

\section{局部多项式逼近}

\begin{frame}{为什么要用 Taylor 展开}
\small
研究函数 $f$ 在 $x_0$ 附近的局部性态时，最自然的想法是用简单函数逼近它。

\begin{block}{逼近层次}
\begin{itemize}
  \item 若 $f$ 在 $x_0$ 连续，则可用常数 $f(x_0)$ 逼近。
  \item 若 $f$ 在 $x_0$ 可微，则可用切线函数逼近。
  \item 若 $f$ 有更高阶导数，则可用更高次多项式逼近。
\end{itemize}
\end{block}

Taylor 展开的核心问题是：怎样选多项式，才能让误差尽可能小？
\end{frame}

\begin{frame}{零阶逼近：连续性}
\small
如果 $f$ 在 $x_0$ 处连续，则
\[
  f(x)-f(x_0)=\ord(1)\qquad (x\to x_0).
\]

这说明在 $x_0$ 附近，
\[
  f(x)\approx f(x_0).
\]

\begin{block}{读法}
这里的 $\ord(1)$ 表示误差趋于 $0$。连续性给出的只是“函数值接近”，还没有给出误差和 $x-x_0$ 的阶数关系。
\end{block}
\end{frame}

\begin{frame}{一阶逼近：切线}
\small
如果 $f$ 在 $x_0$ 处可微，则
\[
  f(x)-\bigl[f(x_0)+f'(x_0)(x-x_0)\bigr]
  =
  \ord(x-x_0).
\]

即在 $x_0$ 附近，$f$ 可由线性函数
\[
  L(x)=f(x_0)+f'(x_0)(x-x_0)
\]
逼近。

\begin{block}{关键理解}
可微的本质不是“有一条切线”，而是函数与切线的误差比 $x-x_0$ 更小。
\end{block}
\end{frame}

\begin{frame}{二阶逼近：抛物线}
\small
如果 $f''(x_0)$ 存在，则
\[
  f(x)-\left[
  f(x_0)+f'(x_0)(x-x_0)
  +\frac12 f''(x_0)(x-x_0)^2
  \right]
  =
  \ord\bigl((x-x_0)^2\bigr).
\]

\begin{block}{从线性到二次}
二阶项描述曲线相对于切线的弯曲程度。若只做极限计算，常常展开到第一个非零项就够；若要证明不等式或估计误差，则需要保留余项信息。
\end{block}
\end{frame}

\begin{frame}{展开前先统一记号}
\small
令
\[
  h=x-x_0.
\]

Taylor 展开常写成
\[
  f(x_0+h)
  =
  f(x_0)+f'(x_0)h+\frac{f''(x_0)}{2!}h^2+\cdots.
\]

\begin{block}{两个问题}
\begin{itemize}
  \item 多项式部分的系数为什么一定是这些导数？
  \item 截断后剩下的余项到底有多大？
\end{itemize}
\end{block}
\end{frame}

\section{Peano 余项}

\begin{frame}{定理 1：带 Peano 余项的 Taylor 公式}
\small
\begin{block}{定理}
设 $f$ 在 $x_0$ 处 $n$ 阶可导，则当 $x\to x_0$ 时，
\[
\begin{aligned}
f(x)
&= f(x_0)+f'(x_0)(x-x_0)
  +\frac{1}{2!}f''(x_0)(x-x_0)^2+\cdots  \\
&\quad
  +\frac{1}{n!}f^{(n)}(x_0)(x-x_0)^n
  +\ord\bigl((x-x_0)^n\bigr).
\end{aligned}
\]
\end{block}

这叫 $f$ 在 $x_0$ 处的 Taylor 公式，最后一项叫 Peano 余项。
\end{frame}

\begin{frame}{Peano 余项的含义}
\small
记
\[
  P_n(x)=\sum_{k=0}^n
  \frac{f^{(k)}(x_0)}{k!}(x-x_0)^k.
\]

则 Peano 余项说的是
\[
  f(x)-P_n(x)=\ord\bigl((x-x_0)^n\bigr).
\]

\begin{block}{怎么使用}
\begin{itemize}
  \item 做极限：它告诉我们高阶小量可以忽略到哪一级。
  \item 做局部判断：它能判断函数在 $x_0$ 附近的符号、单调、极值。
  \item 但它通常不能给出明确的误差上界。
\end{itemize}
\end{block}
\end{frame}

\begin{frame}{证明思路：从 $n$ 到 $n+1$}
\small
证明可用数学归纳法。假设结论对 $n=k$ 成立。若 $f^{(k+1)}(x_0)$ 存在，则对 $f'$ 使用归纳假设：
\[
\begin{aligned}
f'(x)
&= f'(x_0)+f''(x_0)(x-x_0)
  +\frac{1}{2!}f'''(x_0)(x-x_0)^2+\cdots\\
&\quad
  +\frac{1}{k!}f^{(k+1)}(x_0)(x-x_0)^k
  +\ord\bigl((x-x_0)^k\bigr).
\end{aligned}
\]

接下来用洛必达法则，把 $f$ 的余项除以 $(x-x_0)^{k+1}$。
\end{frame}

\begin{frame}{证明思路：余项比值趋于 $0$}
\footnotesize
要证明
\[
\begin{aligned}
f(x)
&= f(x_0)+f'(x_0)(x-x_0)+\cdots\\
&\quad+\frac{1}{(k+1)!}f^{(k+1)}(x_0)(x-x_0)^{k+1}
 +\ord\bigl((x-x_0)^{k+1}\bigr),
\end{aligned}
\]
只需证明
\[
\lim_{x\to x_0}
\frac{
f(x)-\left[f(x_0)+f'(x_0)(x-x_0)+\cdots+
\frac{f^{(k+1)}(x_0)}{(k+1)!}(x-x_0)^{k+1}\right]
}{(x-x_0)^{k+1}}
=0.
\]

对上式用洛必达法则，分子求导后正好变成 $f'$ 的 Peano 余项；由归纳假设，极限为 $0$。
\end{frame}

\begin{frame}{Taylor 余项的记号}
\small
记
\[
  R_n(x)=f(x)-P_n(x)
  =
  f(x)-\sum_{k=0}^n
  \frac{f^{(k)}(x_0)}{k!}(x-x_0)^k.
\]

根据定理 1，
\[
  R_n(x)=\ord\bigl((x-x_0)^n\bigr)\qquad (x\to x_0).
\]

\begin{block}{注意}
Peano 余项只描述 $x\to x_0$ 时的局部阶数。如果要在一个区间上估计误差，通常要用 Lagrange 余项或积分余项。
\end{block}
\end{frame}

\section{带具体余项的 Taylor 公式}

\begin{frame}{为什么还需要具体余项}
\small
Peano 余项告诉我们
\[
  R_n(x)=\ord\bigl((x-x_0)^n\bigr),
\]
但它通常不告诉我们 $R_n(x)$ 的具体大小。

\begin{block}{具体余项能解决的问题}
\begin{itemize}
  \item 判断 Taylor 级数是否收敛到原函数。
  \item 给出数值近似的误差上界。
  \item 推导一些积分、级数和不等式的精确估计。
\end{itemize}
\end{block}
\end{frame}

\begin{frame}{定理 2：Lagrange 与 Cauchy 余项}
\footnotesize
\begin{block}{定理}
设 $f$ 在开区间 $(a,b)$ 中有直到 $n+1$ 阶导数，$x_0,x\in(a,b)$。则存在介于 $x_0$ 与 $x$ 之间的点 $\xi$，使得
\[
  R_n(x)=\frac{1}{(n+1)!}f^{(n+1)}(\xi)(x-x_0)^{n+1}.
\]
这叫 Lagrange 余项。

也存在介于 $x_0$ 与 $x$ 之间的点 $\xi$，使得
\[
  R_n(x)=\frac{1}{n!}f^{(n+1)}(\xi)(x-\xi)^n(x-x_0).
\]
这叫 Cauchy 余项。
\end{block}
\end{frame}

\begin{frame}{辅助函数 $F(t)$}
\small
证明时把 $x$ 看成固定点，以 $t$ 为变量，构造
\[
  F(t)=f(t)+\sum_{k=1}^n
  \frac{f^{(k)}(t)}{k!}(x-t)^k,\qquad t\in(a,b).
\]

对 $t$ 求导，可得大量项相消：
\[
\begin{aligned}
F'(t)
&=f'(t)+\sum_{k=1}^n
\left[
\frac{f^{(k+1)}(t)}{k!}(x-t)^k
-\frac{f^{(k)}(t)}{(k-1)!}(x-t)^{k-1}
\right] \\
&=\frac1{n!}f^{(n+1)}(t)(x-t)^n.
\end{aligned}
\]
\end{frame}

\begin{frame}{从 $F(t)$ 得到 Cauchy 余项}
\small
由 $F$ 的构造，
\[
  F(x)-F(x_0)=R_n(x).
\]

由 Lagrange 微分中值定理，存在介于 $x_0$ 与 $x$ 之间的点 $\xi$，使得
\[
  R_n(x)=F'(\xi)(x-x_0).
\]

代入 $F'(\xi)$，得到
\[
  R_n(x)=\frac1{n!}f^{(n+1)}(\xi)(x-\xi)^n(x-x_0).
\]

这就是 Cauchy 余项。
\end{frame}

\begin{frame}{从 $F(t)$ 得到 Lagrange 余项}
\small
取
\[
  G(t)=-(x-t)^{n+1}.
\]

由 Cauchy 微分中值定理，存在介于 $x_0$ 与 $x$ 之间的点 $\xi$，使得
\[
  \frac{R_n(x)}{(x-x_0)^{n+1}}
  =
  \frac{F(x)-F(x_0)}{G(x)-G(x_0)}
  =
  \frac{F'(\xi)}{G'(\xi)}.
\]

因为
\[
  G'(\xi)=(n+1)(x-\xi)^n,
\]
所以
\[
  R_n(x)=\frac{1}{(n+1)!}f^{(n+1)}(\xi)(x-x_0)^{n+1}.
\]
\end{frame}

\begin{frame}{端点情形的说明}
\small
如果 $f$ 的直到 $n$ 阶导数在 $[a,b]$ 上连续，且 $f$ 在 $(a,b)$ 中有 $n+1$ 阶导数，则对端点 $x=a$ 或 $x=b$，Taylor 余项公式仍然成立。

\begin{block}{原因}
微分中值定理需要的是区间内部可导、闭区间连续。端点处不要求存在所有高阶导数，只要能保证辅助函数在闭区间连续即可。
\end{block}

这点在用 Taylor 公式估计闭区间上的误差时很常见。
\end{frame}

\begin{frame}{引理 1：零点型 Taylor 公式}
\small
\begin{block}{引理}
设 $g$ 在 $(a,b)$ 内 $n+1$ 阶可导，$x_0\in(a,b)$。如果
\[
  g(x_0)=g'(x_0)=\cdots=g^{(n)}(x_0)=0,
\]
则任给 $c\in(a,b)$，存在介于 $x_0$ 与 $c$ 之间的点 $\xi$，使得
\[
  g(c)=\frac{g^{(n+1)}(\xi)}{(n+1)!}(c-x_0)^{n+1}.
\]
\end{block}

这是 Lagrange 余项的一个常用证明工具。
\end{frame}

\begin{frame}{引理 1 的证明思路}
\small
假设 $c\ne x_0$。构造
\[
  h(x)=g(x)-\frac{g(c)}{(c-x_0)^{n+1}}(x-x_0)^{n+1}.
\]

则
\[
  h(c)=0,\qquad h(x_0)=h'(x_0)=\cdots=h^{(n)}(x_0)=0.
\]

反复使用 Rolle 中值定理，可以在 $c$ 与 $x_0$ 之间找到一点 $\xi$，使得
\[
  h^{(n+1)}(\xi)=0.
\]

整理即得
\[
  g(c)=\frac{g^{(n+1)}(\xi)}{(n+1)!}(c-x_0)^{n+1}.
\]
\end{frame}

\begin{frame}{积分余项}
\small
如果 $f$ 的条件更强，例如 $f$ 是 $n+1$ 阶连续可微，则由微积分基本公式可得
\[
  R_n(x)
  =
  \int_{x_0}^{x}
  \frac{f^{(n+1)}(t)}{n!}(x-t)^n\,\dd t.
\]

\begin{block}{积分余项的优点}
\begin{itemize}
  \item 它不是“存在某个 $\xi$”的形式，而是一个精确积分表达式。
  \item 适合处理积分、估计以及某些级数求和问题。
  \item 由积分中值定理还能重新得到 Lagrange 余项。
\end{itemize}
\end{block}
\end{frame}

\begin{frame}{三种余项怎么选}
\small
\begin{center}
\begin{tabular}{c|c|c}
余项类型 & 适合场景 & 主要信息\\
\hline
Peano 余项 & 极限、局部符号 & 阶数足够小\\
Lagrange 余项 & 误差上界、收敛性 & 有明确估计\\
Cauchy 余项 & 特殊结构、证明 & 形式更灵活\\
积分余项 & 积分、级数、精确表达 & 可以直接计算或估计
\end{tabular}
\end{center}

\vspace{0.8em}
实际做题时，先问自己：我只需要局部阶数，还是需要误差大小？
\end{frame}

\section{Taylor 级数与常用展开}

\begin{frame}{Taylor 级数与 Maclaurin 展开}
\small
如果 $f$ 在 $x_0$ 附近无限次可微，则形式级数
\[
  \sum_{n=0}^{\infty}
  \frac{f^{(n)}(x_0)}{n!}(x-x_0)^n
\]
称为 $f$ 在 $x_0$ 处的 Taylor 级数。

当 $x_0=0$ 时，称为 Maclaurin 级数。

\begin{block}{关键条件}
只有当
\[
  \lim_{n\to\infty}R_n(x)=0
\]
时，Taylor 级数才收敛到函数 $f(x)$ 本身。
\end{block}
\end{frame}

\begin{frame}{例 1：几何级数}
\small
\textbf{求} $f(x)=\dfrac1{1-x}$ \textbf{在} $x=0$ \textbf{处的 Taylor 展开，并判断收敛性。}

\vspace{0.6em}
因为
\[
  f^{(n)}(x)=\frac{n!}{(1-x)^{n+1}},
\]
所以
\[
  f^{(n)}(0)=n!.
\]

因此形式展开为
\[
  \frac1{1-x}=1+x+x^2+\cdots+x^n+\cdots.
\]
\end{frame}

\begin{frame}{例 1：余项判断}
\small
余项为
\[
  R_n(x)
  =
  \frac1{1-x}-(1+x+x^2+\cdots+x^n)
  =
  \frac{x^{n+1}}{1-x}.
\]

当 $x\in(-1,1)$ 时，
\[
  R_n(x)\to0.
\]

因此
\[
  \frac1{1-x}=1+x+x^2+\cdots+x^n+\cdots,\qquad x\in(-1,1).
\]
\end{frame}

\begin{frame}{例 2：指数函数}
\small
\textbf{求} $\ee^x$ \textbf{在} $x=0$ \textbf{处的 Taylor 展开。}

\vspace{0.6em}
因为 $f^{(n)}(x)=\ee^x$，所以
\[
  f^{(n)}(0)=1.
\]

因此
\[
  \ee^x
  =
  1+x+\frac{x^2}{2!}+\frac{x^3}{3!}
  +\cdots+\frac{x^n}{n!}+R_n(x).
\]

Lagrange 余项为
\[
  R_n(x)=\frac{\ee^{\theta x}}{(n+1)!}x^{n+1},\qquad 0<\theta<1.
\]
\end{frame}

\begin{frame}{例 2：余项趋于 $0$}
\small
对固定的 $x$，有
\[
  \abs{R_n(x)}
  \le
  \ee^{\abs{x}}\frac{\abs{x}^{n+1}}{(n+1)!}
  \to0.
\]

所以
\[
  \ee^x
  =
  1+x+\frac{x^2}{2!}+\frac{x^3}{3!}
  +\cdots+\frac{x^n}{n!}+\cdots,
  \qquad x\in\R.
\]

特别地，
\[
  \ee=1+1+\frac1{2!}+\frac1{3!}+\cdots+\frac1{n!}+\cdots.
\]
\end{frame}

\begin{frame}{例 2：用展开证明 $e$ 不是整数}
\small
由
\[
  \ee=1+1+\frac1{2!}+\frac1{3!}+\cdots
\]
可得
\[
  2<\ee<1+1+\frac12+\frac1{2\cdot3}+\frac1{3\cdot4}+\cdots=3.
\]

因此 $\ee$ 不是整数。

\begin{block}{说明}
这个估计利用了阶乘分母增长很快。Taylor 展开的价值不只是求近似值，还能给出严格不等式。
\end{block}
\end{frame}

\begin{frame}{例 2：进一步证明 $e$ 是无理数}
\footnotesize
反设
\[
  \ee=\frac pq
\]
为有理数。则
\[
  q!\left(\ee-\left(1+1+\frac1{2!}+\cdots+\frac1{q!}\right)\right)
\]
左边是一个正整数。

另一方面，
\[
\begin{aligned}
0
&<q!\left(\ee-\left(1+1+\frac1{2!}+\cdots+\frac1{q!}\right)\right)\\
&=q!\left(\frac1{(q+1)!}+\frac1{(q+2)!}+\cdots\right)\\
&<\frac1{q+1}+\frac1{(q+1)(q+2)}
  +\frac1{(q+2)(q+3)}+\cdots\\
&<\frac2{q+1}<1,
\end{aligned}
\]
矛盾。因此 $\ee$ 是无理数。
\end{frame}

\begin{frame}{例 3：$\sin x$ 的 Taylor 展开}
\small
因为
\[
  \sin^{(n)}x=\sin\left(x+\frac n2\pi\right),
\]
所以
\[
  \sin^{(2k+1)}(0)=(-1)^k,\qquad \sin^{(2k)}(0)=0.
\]

由 Lagrange 余项公式，
\[
  \sin x
  =
  x-\frac{x^3}{3!}+\frac{x^5}{5!}
  -\cdots+\frac{(-1)^n x^{2n+1}}{(2n+1)!}
  +R_{2n+1}(x).
\]

由于余项趋于 $0$，故
\[
  \sin x=\sum_{n=0}^{\infty}\frac{(-1)^n x^{2n+1}}{(2n+1)!},
  \qquad x\in\R.
\]
\end{frame}

\begin{frame}{例 3：$\cos x$ 的 Taylor 展开}
\small
类似地，
\[
  \cos x
  =
  \sum_{n=0}^{\infty}
  \frac{(-1)^n x^{2n}}{(2n)!}.
\]

也就是
\[
  \cos x
  =
  1-\frac{x^2}{2!}+\frac{x^4}{4!}
  -\frac{x^6}{6!}+\cdots,\qquad x\in\R.
\]

\begin{block}{记忆方式}
\begin{itemize}
  \item $\sin x$ 是奇函数，所以只出现奇次幂。
  \item $\cos x$ 是偶函数，所以只出现偶次幂。
  \item 符号交替来自导数循环。
\end{itemize}
\end{block}
\end{frame}

\begin{frame}{定理 3：Taylor 系数的唯一性}
\small
\begin{block}{定理}
设 $f$ 在 $x_0$ 处 $n$ 阶可导，且
\[
  f(x)=\sum_{k=0}^n a_k(x-x_0)^k
  +\ord\bigl((x-x_0)^n\bigr)\qquad (x\to x_0).
\]
则
\[
  a_k=\frac1{k!}f^{(k)}(x_0),
  \qquad k=0,1,\ldots,n.
\]
\end{block}

这说明：只要展开到同一阶，Taylor 多项式的系数不可能有第二种选择。
\end{frame}

\begin{frame}{唯一性的证明思路}
\small
由 Peano 余项公式，
\[
  f(x)=\sum_{k=0}^n
  \frac1{k!}f^{(k)}(x_0)(x-x_0)^k
  +\ord\bigl((x-x_0)^n\bigr).
\]

两式相减，令
\[
  b_k=a_k-\frac1{k!}f^{(k)}(x_0),
\]
则
\[
  \sum_{k=0}^n b_k(x-x_0)^k
  =
  \ord\bigl((x-x_0)^n\bigr).
\]

依次令 $x\to x_0$、两边除以 $(x-x_0)$ 再令 $x\to x_0$，可得
\[
  b_0=b_1=\cdots=b_n=0.
\]
\end{frame}

\begin{frame}{命题 1：Taylor 级数的运算}
\footnotesize
若
\[
  f(x)=\sum_{n=0}^{\infty}a_nx^n
\]
是 $f$ 在 $0$ 处的 Taylor 展开，则形式上有
\[
  f(-x)=\sum_{n=0}^{\infty}(-1)^n a_nx^n,
\]
\[
  f(x^k)=\sum_{n=0}^{\infty}a_nx^{kn}\qquad (k\in\mathbb{N}^*),
\]
\[
  x^k f(x)=\sum_{n=0}^{\infty}a_nx^{n+k}.
\]

这些结论来自 Taylor 系数的唯一性。
\end{frame}

\begin{frame}{命题 1：求导、积分和线性组合}
\footnotesize
在收敛性允许的范围内，
\[
  f'(x)=\sum_{n=1}^{\infty}na_nx^{n-1}
  =
  \sum_{n=0}^{\infty}(n+1)a_{n+1}x^n,
\]
\[
  \int_0^x f(t)\,\dd t
  =
  \sum_{n=0}^{\infty}\frac{a_n}{n+1}x^{n+1}.
\]

若
\[
  g(x)=\sum_{n=0}^{\infty}b_nx^n,
\]
则
\[
  \lambda f(x)+\mu g(x)
  =
  \sum_{n=0}^{\infty}(\lambda a_n+\mu b_n)x^n.
\]
\end{frame}

\begin{frame}{由几何级数推出的展开}
\small
由
\[
  \frac1{1-x}=1+x+x^2+x^3+\cdots,\qquad x\in(-1,1),
\]
可得
\[
  \frac1{1+x}=1-x+x^2-x^3+x^4-\cdots,\qquad x\in(-1,1),
\]
\[
  \frac1{1-x^2}=1+x^2+x^4+x^6+\cdots,\qquad x\in(-1,1),
\]
\[
  \frac1{1+x^2}=1-x^2+x^4-x^6+\cdots,\qquad x\in(-1,1).
\]
\end{frame}

\begin{frame}{由积分得到 $\ln(1-x)$}
\small
因为
\[
  \frac1{1-t}=1+t+t^2+\cdots+t^{n-1}+\frac{t^n}{1-t},
\]
所以
\[
\begin{aligned}
\ln(1-x)
&=-\int_0^x\frac{\dd t}{1-t}\\
&=-x-\frac{x^2}{2}-\frac{x^3}{3}
  -\cdots-\frac{x^n}{n}
  +R_n(x),
\end{aligned}
\]
其中
\[
  R_n(x)=-\int_0^x\frac{t^n}{1-t}\,\dd t.
\]
\end{frame}

\begin{frame}{余项估计：$\ln(1-x)$}
\small
若 $-1\le x<0$，则
\[
  \abs{R_n(x)}
  \le
  \int_0^{\abs{x}}t^n\,\dd t
  =
  \frac{\abs{x}^{n+1}}{n+1}\to0.
\]

若 $0\le x<1$，则
\[
  \abs{R_n(x)}
  \le
  \frac1{1-x}\int_0^x t^n\,\dd t
  =
  \frac{x^{n+1}}{(1-x)(n+1)}\to0.
\]

因此
\[
  \ln(1-x)=-x-\frac{x^2}{2}-\frac{x^3}{3}-\cdots,\qquad x\in[-1,1).
\]
\end{frame}

\begin{frame}{由换元得到 $\ln(1+x)$}
\small
在
\[
  \ln(1-x)=-x-\frac{x^2}{2}-\frac{x^3}{3}-\cdots
\]
中令 $x$ 换成 $-x$，得到
\[
  \ln(1+x)
  =
  x-\frac{x^2}{2}+\frac{x^3}{3}
  -\frac{x^4}{4}+\frac{x^5}{5}-\cdots.
\]

成立范围为
\[
  x\in(-1,1].
\]

特别地，取 $x=1$，得到
\[
  \ln2=1-\frac12+\frac13-\frac14+\frac15-\frac16+\cdots.
\]
\end{frame}

\begin{frame}{由积分得到 $\arctan x$}
\small
因为
\[
  \arctan x=\int_0^x\frac{\dd t}{1+t^2},
\]
且
\[
  \frac1{1+t^2}
  =
  1-t^2+t^4-\cdots+(-1)^{n-1}t^{2n-2}
  +\frac{(-1)^nt^{2n}}{1+t^2},
\]
所以
\[
  \arctan x
  =
  x-\frac{x^3}{3}+\frac{x^5}{5}
  -\cdots+(-1)^{n-1}\frac{x^{2n-1}}{2n-1}
  +R_n(x).
\]
\end{frame}

\begin{frame}{余项估计：$\arctan x$}
\small
其中
\[
  R_n(x)=(-1)^n\int_0^x\frac{t^{2n}}{1+t^2}\,\dd t.
\]

当 $x\in[-1,1]$ 时，
\[
  \abs{R_n(x)}
  \le
  \int_0^{\abs{x}}t^{2n}\,\dd t
  =
  \frac{\abs{x}^{2n+1}}{2n+1}\to0.
\]

因此
\[
  \arctan x=x-\frac{x^3}{3}+\frac{x^5}{5}-\frac{x^7}{7}+\cdots,
  \qquad x\in[-1,1].
\]

取 $x=1$，得
\[
  \frac\pi4=1-\frac13+\frac15-\frac17+\cdots.
\]
\end{frame}

\begin{frame}{更多常用展开}
\small
由已经得到的展开式，还可推出
\[
  \ee^{-x}
  =
  1-x+\frac{x^2}{2!}-\frac{x^3}{3!}
  +\frac{x^4}{4!}-\cdots
  =
  \sum_{n=0}^{\infty}(-1)^n\frac{x^n}{n!},
\]
\[
  \sinh x
  =
  \frac12(\ee^x-\ee^{-x})
  =
  x+\frac{x^3}{3!}+\frac{x^5}{5!}+\cdots
  =
  \sum_{n=0}^{\infty}\frac{x^{2n+1}}{(2n+1)!},
\]
\[
  \cosh x
  =
  \frac12(\ee^x+\ee^{-x})
  =
  \sum_{n=0}^{\infty}\frac{x^{2n}}{(2n)!}.
\]
\end{frame}

\begin{frame}{常用展开表}
\footnotesize
\[
\begin{array}{c|c}
\text{函数} & \text{Maclaurin 展开}\\
\hline
\ee^x &
1+x+\dfrac{x^2}{2!}+\dfrac{x^3}{3!}+\cdots\\[0.6em]
\sin x &
x-\dfrac{x^3}{3!}+\dfrac{x^5}{5!}-\cdots\\[0.6em]
\cos x &
1-\dfrac{x^2}{2!}+\dfrac{x^4}{4!}-\cdots\\[0.6em]
\ln(1+x) &
x-\dfrac{x^2}{2}+\dfrac{x^3}{3}-\cdots\\[0.6em]
\dfrac1{1-x} &
1+x+x^2+x^3+\cdots\\[0.6em]
\arctan x &
x-\dfrac{x^3}{3}+\dfrac{x^5}{5}-\cdots
\end{array}
\]

使用时必须同时记住展开中心和适用范围。
\end{frame}

\section{积分余项的应用}

\begin{frame}{例 4：积分余项的一个应用}
\small
\textbf{求}
\[
  \int_0^1(1-t^2)^n\,\dd t.
\]

考虑多项式函数
\[
  f(x)=(1+x)^{2n+1}.
\]

在 $x_0=0$ 处展开到 $n$ 阶：
\[
  (1+x)^{2n+1}
  =
  \sum_{k=0}^n\binom{2n+1}{k}x^k+R_n(x).
\]
\end{frame}

\begin{frame}{例 4：写出积分余项}
\small
因为
\[
  f^{(n+1)}(t)=\frac{(2n+1)!}{n!}(1+t)^n,
\]
所以
\[
  R_n(x)
  =
  \frac1{n!}\int_0^x
  \frac{(2n+1)!}{n!}(1+t)^n(x-t)^n\,\dd t.
\]

特别地，取 $x=1$，得到
\[
  R_n(1)=
  \frac{(2n+1)!}{(n!)^2}
  \int_0^1(1-t^2)^n\,\dd t.
\]
\end{frame}

\begin{frame}{例 4：完成计算}
\small
另一方面，
\[
  R_n(1)
  =
  2^{2n+1}-\sum_{k=0}^n\binom{2n+1}{k}.
\]

由二项式系数对称性，
\[
  \sum_{k=0}^n\binom{2n+1}{k}=2^{2n}.
\]

所以
\[
  R_n(1)=2^{2n}.
\]

于是
\[
  \int_0^1(1-t^2)^n\,\dd t
  =
  \frac{(n!)^2}{(2n+1)!}2^{2n}
  =
  \frac{(2n)!!}{(2n+1)!!}.
\]
\end{frame}

\section{收敛但不收敛到函数}

\begin{frame}{一个重要提醒}
\small
很多函数可以得到 Taylor 级数，但要判断它是否收敛、是否收敛到原函数，并不总是容易。

\begin{block}{关键区别}
\begin{itemize}
  \item Taylor 级数是由函数在某一点的所有导数决定的。
  \item 函数在一点的所有导数，不一定决定该点附近的函数值。
  \item 因此存在 Taylor 级数收敛，但不收敛到函数本身的例子。
\end{itemize}
\end{block}
\end{frame}

\begin{frame}{例 5：光滑但 Taylor 展开失效}
\small
定义函数
\[
\varphi(x)=
\begin{cases}
0,&x\le0,\\
\ee^{-1/x},&x>0.
\end{cases}
\]

显然，在 $(-\infty,0)\cup(0,\infty)$ 中，$\varphi$ 是无限次可微的。

进一步可以证明，$\varphi$ 在 $x=0$ 处也无限次可微，且
\[
  \varphi^{(n)}(0)=0,\qquad n\ge0.
\]

因此 $\varphi$ 在 $0$ 处的 Taylor 级数恒为 $0$。
\end{frame}

\begin{frame}{例 5：为什么这是反例}
\small
虽然 $\varphi$ 在 $0$ 处的 Taylor 级数是
\[
  0+0x+0x^2+\cdots=0,
\]
但当 $x>0$ 时，
\[
  \varphi(x)=\ee^{-1/x}>0.
\]

所以 Taylor 级数没有收敛到函数本身。

\begin{block}{结论}
无限可微不等于能够被 Taylor 级数表示。要让 Taylor 级数代表原函数，必须证明余项 $R_n(x)\to0$。
\end{block}
\end{frame}

\section{使用策略}

\begin{frame}{做 Taylor 展开题的基本流程}
\small
\begin{enumerate}
  \item \textbf{确定展开中心：}通常是极限点、近似点或题目指定点。
  \item \textbf{确定展开阶数：}展开到第一个不被抵消的非零项，或题目所需误差阶。
  \item \textbf{选择余项：}只做局部极限用 Peano；要误差估计用 Lagrange 或积分余项。
  \item \textbf{检查范围：}级数展开必须说明收敛区间；有限阶展开要说明 $x\to x_0$。
\end{enumerate}
\end{frame}

\begin{frame}{常见错误}
\small
\begin{itemize}
  \item 只写展开式，不写适用条件。
  \item 忽略余项，导致把局部近似当成全局等式。
  \item 展开阶数不够，抵消后无法判断主项。
  \item 把“Taylor 级数存在”误认为“一定等于原函数”。
\end{itemize}

\begin{block}{最实用的原则}
如果题目中出现抵消，就至少展开到抵消后的第一项；如果题目要求误差或收敛，就必须处理余项。
\end{block}
\end{frame}

\begin{frame}{课堂小结}
\begin{itemize}
  \item Taylor 展开是用多项式描述函数在一点附近的局部性态。
  \item Peano 余项适合极限和局部判断，Lagrange 与积分余项适合误差估计。
  \item 常用展开式要同时记住展开中心、收敛范围和余项控制。
  \item 无限可微不保证 Taylor 级数收敛到函数本身，必须检查 $R_n(x)\to0$。
\end{itemize}
\end{frame}

\end{document}
